% 摘要：100--400 字。問題 → 方法 → 主要結果（具體數字）→ 意義。數字待實驗後填入（紅字 [TBD]）。
Tensor parallelism (TP) is the default way to spread large language model (LLM) inference
across GPUs, but its benefit depends on how fast the GPUs can all-reduce activations.
Datacenter GPUs connect through NVLink, while commodity GPUs exchange data over PCIe,
and practitioners have little guidance on when adding GPUs actually pays off.
We study the vLLM inference engine on two platforms: NVIDIA H100 GPUs connected by NVLink
and RTX~4090 GPUs connected by PCIe.
We measure strong and weak scaling of TP for the prefill and decode phases, profile where
GPU time goes, and ablate two communication-related engine features: vLLM's custom
all-reduce kernel and CUDA graphs.
We then build a simple analytical model that combines a roofline-style compute and memory
term with an all-reduce term calibrated by a communication microbenchmark.
On H100, TP=2 reaches a decode speedup of \tbd (\tbd\% efficiency).
On RTX~4090, TP=2 and TP=4 reach \tbd and \tbd, and all-reduce accounts for up to
\tbd\% of GPU kernel time.
Calibrated only on single-GPU runs, the model predicts multi-GPU step time with a mean
error of \tbd\%.
We release all scripts and data, designed to run under strict cluster limits of
6-minute, 2-GPU jobs.
